\documentclass[openright, twoside, a4paper]{memoir}
\usepackage[brazil]{babel}
\usepackage{livros_crp}
\usepackage{longtable}
\title{Política de Segurança em Tecnologia da Informação e Comunicação
(PSTIC) do Conselho Regional de Psicologia da 6ª Região}
\author{Unidade de TIC — CRP-06}
\date{2024}
\begin{document}
\frontmatter
%

\chapter[Apresentação da 2ª edição]{Apresentação da 2ª edição}

Passados dez meses de implantação da atual Política de Segurança em
Tecnologia da Informação e Comunicação (PSTIC), novos recursos e fluxos
tecnológicos foram implantados, tornando necessária sua atualização.

O desafio de instituir uma política institucional dentro de uma
autarquia pública federal que, durante muitos anos, atuou em processos
diversos e difusos sem padronização e formalidade, coloca-nos em
situação de maior atenção e responsabilidade.

Identificamos resistências e dúvidas quanto aos procedimentos
institucionais e maior necessidade de elucidação das ações.

O desafio de garantir a governança preconizada para a administração
pública envolve três funções básicas, alinhadas às tarefas sugeridas
pela ISO/IEC 38500:2008:

\begin{enumerate}
\item avaliar o ambiente, os cenários, o desempenho e os resultados atuais
    e futuros;

\item direcionar e orientar a preparação, a articulação e a coordenação de
    políticas e planos, alinhando as funções organizacionais às necessidades
    das partes interessadas (usuária/os dos serviços, cidadãos e sociedade
    em geral) e assegurando o alcance dos objetivos estabelecidos;

\item monitorar os resultados, o desempenho e o cumprimento de políticas e
    planos, confrontando-os com as metas estabelecidas e as expectativas das
    partes interessadas.
\end{enumerate}

O Acórdão 1768/2022 do TCU visa contribuir para a transformação digital
do país, conscientizando os gestores públicos acerca dos riscos aos
quais as organizações estão sujeitas, de modo que sejam implementados
controles e medidas de segurança adequados para enfrentá-los.

Assim, o TCU estabelece os ``Cinco controles de segurança cibernética
para ontem'' para embasar a adoção, pelas organizações públicas
federais, dos controles críticos de cibersegurança que se seguem:

\begin{itemize}[]
\item \textbf{Controle 1 --- Inventário e controle de ativos corporativos}:
    identificar e impedir a utilização de ativos de TI não
    autorizados/gerenciados como vetores de ataques cibernéticos;

\item \textbf{Controle 2 --- Inventário e controle de ativos de
      \emph{software}}: identificar e impedir a utilização de \emph{softwares}
    não autorizados/gerenciados como vetores de ataques cibernéticos;

\item \textbf{Controle 7 --- Gestão contínua de vulnerabilidades}:
    evitar a exploração de vulnerabilidades conhecidas nos ativos
    corporativos de TI;

\item \textbf{Controle 14 --- Conscientização sobre segurança e
      treinamento de competências}: reduzir a possibilidade de incidentes e
    ataques derivados do comportamento humano (engenharia social);

\item \textbf{Controle 17 --- Gestão de respostas a incidentes}: melhorar a
    capacidade de identificar potenciais ameaças e ataques, evitar que se
    espalhem e recuperar rapidamente dados e sistemas eventualmente
    corrompidos.
\end{itemize}

Também a publicação da Portaria CRP-06 nº~99/2024, que ``estabelece a
jornada de trabalho de 30 horas semanais e o teletrabalho no âmbito do
Conselho Regional de Psicologia da 6ª Região, cria e regulamenta o
regime híbrido de trabalho e dá outras providências'' nos coloca a
necessidade de regulamentar com mais rigor o uso das tecnologias da
informação e comunicação.

Esta nova edição apresenta inovações em diferentes âmbitos, seja pela
inserção de tutoriais de acesso aos sistemas e programas, seja pela
atualização de ferramentas para abertura de chamados e reformulação de
fluxos;

propõe uma ponte com o Código de Ética da Sociedade Brasileira de
Computação, um avanço na vinculação prática do compromisso ético de
nossa profissão aos princípios ali indicados para a oferta de serviços
de qualidade a sociedades;

incorpora normativas convergentes na defesa intransigente da democracia
e da Psicologia: a Lei de Acesso à Informação (LAI; Lei nº~12.527/2011),
o Marco Civil da Internet (Lei nº~12.965/2014) e a Lei Geral de Proteção
de Dados (LGPD; Lei nº~13.709/2018), fundamentando, assim a expressão
prática do espírito destas.

Assim, apresentamos a segunda edição da presente Política com o foco no
aprimoramento, na qualificação e na inovação dos fluxos e procedimentos
institucionais.

Talita Fabiano de Carvalho

Conselheira presidenta do CRP-06
\chapter{Apresentação da 1ª edição}

O Conselho Regional de Psicologia de São Paulo --- 6ª Região (CRP~SP),
entidade dotada de personalidade jurídica de direito público, com
autonomia administrativa e financeira, nos termos da Lei Federal
nº~5.766, de 20 de dezembro de 1971, tem como finalidade fiscalizar o
exercício da profissão de psicóloga/o, competindo-lhe orientar,
disciplinar e zelar pela fiel observância dos princípios
ético-profissionais e contribuir para o desenvolvimento da Psicologia
enquanto ciência e profissão.

Este CRP~SP tem sede na cidade de São Paulo e jurisdição no estado de
São Paulo, possuindo subsedes nas regiões de Alto Tietê (Mogi das
Cruzes), Assis (Assis), Baixada Santista e Vale do Ribeira (Santos),
Bauru (Bauru), Campinas (Campinas), Grande ABC (Santo André),
Metropolitana (São Paulo), Ribeirão Preto (Ribeirão Preto), São José do
Rio Preto (São José do Rio Preto), Sorocaba (Sorocaba) e Vale do Paraíba
e Litoral Norte (Taubaté).

O CRP~SP tem como atribuições principais: adotar as medidas e
procedimentos necessários à permanente orientação, disciplina e
fiscalização do exercício da profissão de psicóloga/o no estado de São
Paulo; adotar medidas e procedimentos para preservação do livre
exercício da profissão de psicóloga/o, bem como do respeito às suas
prerrogativas e direitos profissionais; executar os serviços
concernentes ao registro profissional das/os psicólogas/os, realizando
as inscrições, reativações, transferências e cancelamentos de registros,
expedindo às/aos inscritas/os a carteira de identidade profissional
(CIP); funcionar como tribunal regional de ética profissional; elaborar
e encaminhar ao Conselho Federal de Psicologia (CFP) a proposta
orçamentária anual, além do relatório geral de suas atividades;
providenciar a instalação da Assembleia Geral das/os psicólogas/os
inscritas/os na região; arrecadar anuidades, taxas e demais rendimentos
que lhe competem, promovendo o repasse da arrecadação ao Conselho
Federal; criar os serviços necessários ao bom desempenho de suas
atividades e autorizar a compra de material para suas instalações;
organizar o quadro de pessoal; zelar pela dignidade, independência,
prerrogativas e valorização da Psicologia; e zelar pela gestão
responsável, cumprindo a legislação a partir dos princípios da
administração pública: legalidade, impessoalidade, moralidade,
publicidade e eficiência.

Assim, para uma adequada utilização dos recursos e dispositivos
públicos, foi aprovado, no Planejamento Estratégico 2023--2025,
resultado específico para garantir a implantação adequada dos sistemas e
instrumentos oficiais do CRP~SP.

O Resultado 1.2 do Planejamento Estratégico supracitado se refere a
``ter implementado estrutura de gestão democrática com processos de
trabalho planejados e institucionalizados, de forma transversal,
acessível, integrada, transparente e com produção de dados.''

Das ações previstas para 2023, foi atingido o Resultado 1.2: ``Ter
implantado e integrado os sistemas informatizados (SEI, BRC, Benner e
produtos Zimbra/Implanta), garantindo o cumprimento das normativas
(leis, resoluções, Corep), dos fluxos e processos de trabalho com
formação permanente, transparência, produção, análise e segurança de
dados.''

E, por fim, fazendo cumprir a Macroação nº~8 --- ``Ter elaborado
política que institucionaliza e da segurança ao uso dos novos sistemas''
---, é elaborada esta Política de Segurança em Tecnologia da Informação
e Comunicação (PSTIC), que especifica os controles internos aplicáveis à
segurança, ao sigilo da informação e ao uso das ferramentas do Conselho
Regional de Psicologia da 6ª Região --- CRP~SP, com o objetivo de
realizar suas operações gerais, ainda que em situações adversas.

Estão sujeitas/os ao disposto no presente documento todas/os as/os
conselheiros, colaboradoras/es, trabalhadoras/es, funcionárias/os
terceirizadas/os, estagiárias/os, prestadoras/es de serviços e demais
convidadas/os, no que a cada um for aplicável.

A Política Segurança da TIC é o documento que normatiza, orienta e
estabelece as diretrizes para a proteção dos sistemas, documentos,
operações e comunicação institucional quanto a processos de trabalho e
uso da rede e, principalmente com a prevenção de responsabilidade legal
para todos as/os usuárias/os.

Além disso estabelecimento de fluxos para comunicação e guarda de
documentos e deve, portanto, ser cumprida e aplicada em todas as áreas
da autarquia.

A presente política está baseada nas leis vigentes brasileiras que
tangem sobre isso, bem como nas recomendações propostas pelas ABNT NBR
ISO/IEC 27002/2022, reconhecida mundialmente como um código de prática
para a gestão da segurança da informação.

Também é importante destacar que esta política visa apoiar inicialmente
a Lei Geral de Proteção dos Dados --- LGPD (Lei nº~13.709/2018) no qual
toda a informação produzida ou recebida como resultado da atividade
institucional, pertence ao CRP~SP.

Por fim, garantimos aqui que os registros institucionais sejam
disponibilizados no Portal da Transparência do CRP~SP preconizados pela
Lei de Acesso à Informação (Lei nº~12.527/2011) tomadas as devidas
proteções ao sigilo daqueles documentos internos que ficarão guardados
pela Autarquia.

Talita Fabiano de Carvalho

Conselheira presidenta do CRP-06
%
\tableofcontents*
%
\mainmatter
%
\chapter{Dos objetivos e responsabilidades}

\section{Dos objetivos}

\begin{enumerate}

\item
    Estabelecer diretrizes que permitam a todas as pessoas usuárias dos
    sistemas e ferramentas de tecnologia da informação e comunicação (TIC)
    do CRP--06 seguir padrões de comportamento relacionados à segurança da
    informação adequados às necessidades institucionais e à proteção legal
    das/os envolvidas/os;
\item
    normatizar e orientar todas as pessoas usuárias para a proteção dos
    sistemas, documentos, operações e comunicação institucional no que
    tange aos processos de trabalho e uso das redes, principalmente com a
    prevenção de responsabilidade legal;
\item
    definir normas e procedimentos específicos de segurança da informação,
    bem como a implementação de controles e processos para seu
    atendimento.
\item
    ofertar orientações sobre a utilização dos recursos computacionais, de
    telecomunicação e de infraestrutura de TI;
\item
    garantir a integridade da informação para que seja mantida em seu
    estado original, na guarda ou transmissão, contra alterações
    indevidas, intencionais ou acidentais;
\item
    garantir a confidencialidade e que o acesso à informação seja obtido
    somente por pessoas autorizadas e exclusivamente em função do
    interesse público;
\item
    garantir a disponibilidade para que as/os usuárias/os autorizadas/os
    obtenham acesso à informação e aos ativos correspondentes sempre que
    necessário;
\item
    assegurar a aplicabilidade das referências e diretrizes da governança
    da informação preconizadas pelo Tribunal de Contas da União (TCU);
\item
    organizar fluxos institucionais para tramitação de processos,
    informações, solicitações e respostas;
\item
    garantir a lisura, transparência e responsabilidade para com os
    processos institucionais.
\end{enumerate}

\section{Das responsabilidades}

Os recursos institucionais disponíveis na autarquia são de uso exclusivo
da e para as atividades e ações institucionais.

\begin{enumerate}

\item
    O CRP-06 não se responsabiliza pelas perdas causadas pelo uso
    indevido, negligente ou imprudente dos recursos e/ou serviços
    concedidos a suas/seus usuárias/os;
\end{enumerate}

\begin{enumerate}
\item
    o CRP-06, por meio da Unidade de Tecnologia da Informação e
    Comunicação, poderá registrar todo o uso dos sistemas e serviços,
    visando garantir a disponibilidade e a segurança das informações
    utilizadas, inclusive revisando e atualizando estas normas sempre que
    algum fato e/ou evento relevante acontecer;
\item
    todos os contratos firmados com as/os usuárias/os do CRP-06 deverão
    ser previamente formalizados via termo de acordo do uso e
    confidencialidade (TAUC), condição imprescindível para que possa ser
    concedido o acesso aos ativos de informação disponibilizados pela
    instituição;
\item
    qualquer incidente que afete a segurança da informação deverá ser
    comunicado inicialmente à coordenação de TIC e ela, após análise
    encaminhará à Gerência Administrativa e de Tecnologia da Informação
    (GATI) para que as medidas corretivas possam ser tomadas. Se
    necessário, a Unidade de Gestão de Pessoas e\emph{/}ou a diretoria
    poderão ser envolvidas para apoiar a decisão;
\item
    a fim de garantir maior segurança no uso dos sistemas da autarquia,
    faz-se necessária a segregação dos ambientes de desenvolvimento (uso
    da TI), testes (uso da TI e colaboradores-chave para testes),
    homologação (uso da TI) e produção (uso institucional);
\item
    o CRP-06 reserva-se o direito de analisar dados e evidências para
    obtenção de provas a serem utilizadas nos processos investigatórios,
    bem como adotar as medidas legais cabíveis;
\item
    é vedado o uso de todo e qualquer meio não institucional de contas e
    dispositivos pessoais como \emph{e-mails} que não sejam
    institucionais, WhatsApp, Skype, redes sociais e outros fora do
    domínio da rede formal;
\item
    é importante destacar que nenhuma pessoa deve ser lesada no uso de
    seus recursos pessoais, pois a autarquia oferece ferramentas
    necessárias para a comunicação institucional;
\item
    qualquer outro instrumento institucional (\emph{e-mail}, \emph{chat},
    aplicativo de mensagem instantânea, entre outros) não possui a
    prerrogativa deliberativa nem substitui, sob qualquer hipótese,
    aqueles já previstos, como ofício, memorando, despacho, portaria,
    instrução normativa ou resolução;
\item
    nas situações de não cumprimento, parcial ou integral, por parte da/o
    usuária/o, dos requisitos previstos nesta Política e das orientações
    sobre a segurança da informação, bem como das regras internas da
    autarquia, a/o usuária/o será submetida/o às medidas administrativas e
    legais cabíveis;
\item
    o uso pessoal dos recursos institucionais em hipótese alguma possui
    caráter deliberativo ou legal e é expressamente vedado e, caso ocorra,
    resultará na aplicação de medidas previstas na legislação pertinente
    para apuração dos fatos e responsabilização das/os envolvidas/os;
\item
    será de inteira responsabilidade de cada usuária/o o prejuízo ou dano
    que vier a sofrer ou causar ao CRP-06 e/ou a terceiros, em decorrência
    do não cumprimento às diretrizes e normas aqui referidas.
\end{enumerate}

\subsection{Das responsabilidades
específicas}

\subsubsection{Das/os conselheiras/os, gestoras/es, funcionárias/os,
colaboradoras/es, terceirizadas/os, estagiárias/os, prestadoras/es de
serviços e demais
convidadas/os}

\begin{enumerate}

\item
    Zelar pelo cumprimento desta Política, de forma a garantir a segurança
    da informação e assegurar o uso adequado dos meios, realizando
    \emph{backups} periódicos dentro da infraestrutura disponível, não
    compartilhando senhas e não disseminando notícias ou informações
    inconsistentes, sigilosas ou institucionais fora dos instrumentos da
    autarquia;
\end{enumerate}

\begin{enumerate}

\setcounter{enumi}{21}
\item
    exigir das/os usuárias/os sob sua responsabilidade o devido
    cumprimento desta Política;
\item
    exigir e condicionar o uso à assinatura do termo de uso e
    confidencialidade e o seguimento das normas estabelecidas, bem como ao
    compromisso de manter sigilo e confidencialidade sobre todos os ativos
    de informações do CRP-06, elucidando dúvidas e encaminhando para
    suporte as dificuldades técnicas, quando necessário.
\end{enumerate}

\subsubsection{Da Unidade Administrativa de Tecnologia da
Informação}

\begin{enumerate}

\item
    Informar às/aos gestores os serviços específicos que serão prestados e
    os procedimentos de resposta aos incidentes, atualizando, assim, o
    manual de serviços sempre que necessário;
\end{enumerate}

\begin{enumerate}

%\setcounter{enumi}{23}
\item
    configurar os equipamentos, ferramentas e sistemas concedidos às/os
    usuárias/os com todos os controles necessários para cumprir e
    assegurar os requisitos de segurança e uso estabelecidos por esta
    Política;
\item
    permitir, quando necessária, a execução de atividades operacionais sob
    sua responsabilidade, como, por exemplo, manutenção de computadores,
    realização de cópias de segurança, criação de contas para acesso
    institucional ou testes no ambiente;
\item
    seguir as normas e/ou orientações do Sistema Conselhos de Psicologia e
    da legislação vigente, no que diz respeito à segurança da informação
    para todos os Sistemas e aplicações com acesso público;
\item
    planejar, implantar, fornecer e monitorar a capacidade de armazenagem,
    processamento e transmissão necessária para garantir a segurança
    requerida pelas áreas de atuação, fundamentando tecnicamente seus
    procedimentos e encaminhamentos;
\item
    atribuir cada conta ou dispositivo de acesso (computadores, celulares,
    \emph{tablets} e outros), dispositivos dos sistemas, bases de dados ou
    qualquer outro ativo de informação a uma/um responsável identificável
    como pessoa física, de acordo com a Lei Geral de Proteção de Dados
    (LGPD);
\item
    proteger continuamente todos os ativos de informação do CRP-06 contra
    código malicioso, garantindo que os novos ativos só entrem para o
    ambiente de produção após estarem livres de código malicioso e/ou
    indesejado e após auditados por esta unidade administrativa;
\item
    propor a definição de regras formais para instalação de
    \emph{software} e \emph{hardware} em ambiente de produção
    institucional, implementando as diretrizes do Código de ética da
    Sociedade Brasileira de Computação em consonância com o Código de
    Ética Profissional da/o Psicóloga/o;
\item
    realizar manutenções corretiva, evolutiva e preventiva mensalmente, a
    fim de revisar tecnicamente os dispositivos tecnológicos de trabalho e
    mitigar possíveis riscos para assegurar o bom funcionamento dos
    recursos, a partir de plano de trabalho e cronograma devidamente
    aprovados pelo plenário;
\item
    auxiliar na instalação e configuração de assinaturas digitais e
    certificados digitais, por meio de ferramentas específicas;
\item
    promover os meios e estabelecer as diretrizes de orientação para as/os
    usuárias/os quanto à política de \emph{backup} de dados, e auditá-los
    quando necessário;
\item
    garantir, de maneira rápida e eficaz, após solicitação formal, o
    bloqueio de acesso de usuária/o por motivo de desligamento, incidente,
    investigação ou outra situação que exija medida restritiva para fins
    de resguardo dos ativos do CRP-06;
\item
    dar suporte operacional atitudinal na implementação, manutenção e
    promoção da cultura e desenvolvimento das ações e processos para uso
    de todos os sistemas, redes e ativos da autarquia, com base nos
    princípios éticos da área, de modo coerente com esta política.
\end{enumerate}

\paragraph{Da responsabilidade de garantia da segurança de
informação e
comunicação}

\begin{enumerate}

\item
    Publicar e promover as edições necessárias desta Política de Segurança
    em Tecnologia da Informação e Comunicação, com base em princípios
    técnicos;
\item
    promover a conscientização das/os usuárias/os em relação à relevância
    da segurança da informação, mediante campanhas e informativos, gestão
    de termos de cessão de uso e responsabilização;
\item
    analisar criticamente incidentes recorrentes e outros problemas em
    conjunto com usuária/o, fornecedores e, no caso dos Sistemas Federais,
    junto ao Comitê de TI do Conselho Federal de Psicologia;
\item
    analisar técnica e criticamente incidentes recorrentes e outros
    problemas de fornecedoras/es prestadoras/es de serviços em conjunto
    com as gerências e coordenações de cada unidade, elaborar parecer
    técnico fundamentado para dar subsídio à gerência de competência e dar
    ciência ao plenário, efetuando a gestão de contratos de uso e
    prestação de serviços;
\item
    exercer a continua práxis de alinhamento com as diretrizes
    institucionais da autarquia e normas regulamentadoras de segurança da
    informação e comunicação a partir de leis e normativas dos órgãos de
    controle.
\end{enumerate}
\chapter{Da gravação das
reuniões}

Todas as reuniões do CRP~SP devem ser gravadas para fins de registro
posterior, não sendo permitido o compartilhamento das gravações com
terceiros ou com as partes sem autorização expressa de todos os
participantes.

Quanto à proteção dessas gravações, fica determinado que:

\begin{enumerate}

\item
    somente a Unidade de Secretaria terá acesso às gravações para registro
    das reuniões, sendo obrigatória a exclusão do arquivo após a aprovação
    do documento;
\item
    em caso de ata ou registro aprovado ao final da própria reunião, a
    gravação deve ser imediatamente excluída;
\item
    em casos de gravação para fins de treinamento, capacitação e demais
    ações pedagógicas, todas/os as/os participantes devem autorizar
    expressamente sua divulgação no início da atividade, em conjunto com o
    registro da finalidade e da limitação de tal compartilhamento;
\item
    os assuntos discutidos nas reuniões ordinárias, extraordinárias,
    assembleias ou similares e os registros de suas informações são
    passíveis de divulgação e de acesso, com exceção das partes que
    expressamente tiverem sigilo previsto em hipóteses legais e/ou possam
    violar os princípios da Lei Geral de Proteção de Dados Pessoais (LGPD;
    Lei nº~13.709/2018).
\end{enumerate}
\chapter{Do monitoramento e da
auditoria}

Para promover a proteção, é necessário que haja processos básicos, porém
definidos e aplicáveis, delineados em uma estratégia de segurança eficaz
que contemple os três fatores bem estruturados --- pessoas, processos e
tecnologias ---, para que se realize a jornada completa de segurança da
autarquia, mitigando os riscos a que está exposta.

Esses processos visam garantir a integridade, a confidencialidade e a
disponibilidade de dados armazenados e evitar prejuízos. Implementá-los,
bem como zelar pelo rigoroso cumprimento das diretrizes e fluxos aqui
definidos, são algumas das funções da auditoria de segurança da
informação.

A auditoria, portanto, é uma avaliação sistemática e abrangente dos
sistemas de segurança da informação, permitindo sua avaliação e também a
dos processos e políticas de segurança. Seu objetivo principal é
identificar possíveis vulnerabilidades e brechas de segurança para que
medidas de proteção adequadas possam ser implementadas com a frequência
cabível.

Esta Política determina dois tipos de auditoria de segurança da
informação: a interna e a externa.

A auditoria interna, feita pelo Grupo Gestor da TI, é responsável por
verificar e analisar os sistemas e procedimentos internos da
instituição.

A auditoria externa pode ser feita por uma empresa terceirizada e que
não tenha vínculo com a contratante, podendo ser contratada sempre que
necessário.

Essa auditoria visa atingir:

\begin{enumerate}

\item
    garantia de ética e conformidade dos sistemas e dispositivos
    institucionais;
\item
    conformidade com as leis e regulamentações de proteção de dados;
\item
    prevenção de invasões e corrompimento de dados;
\item
    segurança do ambiente de trabalho e confiabilidade;
\item
    garantia de melhor desempenho e produtividade com os dispositivos
    institucionais;
\item
    identificação de riscos de segurança e das vulnerabilidades dos
    sistemas;
\item
    prevenção de futuros problemas;
\item
    identificação de áreas que precisam de atenção e processos mais
    vulneráveis;
\item
    prevenção de perdas financeiras, uso inadequado dos recursos, proteção
    da reputação da autarquia e vazamentos de dados;
\item
    preparação para lidar com ameaças de segurança e recuperar as
    capacidades em caso de falha no sistema ou vazamento de dados.
\end{enumerate}

Assim, realizar auditorias regulares pode ajudar a identificar pontos
fracos na infraestrutura de TI, verificar os controles de segurança e
monitorar o cumprimento destas diretrizes.

Por fim, fica instituída a realização, a qualquer tempo, de inspeção
física ou lógica nos recursos tecnológicos de sua corresponsabilidade,
em caráter de manutenção preventiva, corretiva ou evolutiva, além da
instalação de sistemas de proteção, preventivos e detectáveis, sempre
que possível.

\section{Etapas}

\subsection{Planejamento}

Deverá ser definido anualmente o escopo da auditoria, identificando-se
áreas críticas que precisem ser avaliadas e designando-se uma equipe
responsável pela auditoria. Deverão, também, ser estabelecidos objetivos
e metas para a auditoria e um cronograma para sua execução.

\subsection{Coleta de
informações}

Coletar informações sobre a autarquia, sistemas e processos, o que pode
incluir entrevistas com funcionários, revisão de documentos e análise de
sistemas e infraestrutura.

\subsection{Análise de riscos}

Com base nas informações coletadas, deve-se realizar uma análise de
risco para identificar as principais ameaças e vulnerabilidades da
organização.

\subsection{Avaliação de
controles}

Avaliar os controles de segurança existentes na organização, como
políticas, procedimentos e tecnologias de segurança.

\subsection{Relatórios e
recomendações}

Preparar um relatório detalhado com principais descobertas e
recomendações para mitigar os riscos identificados.

\subsection{Acompanhamento}

Realizar um acompanhamento para garantir que as recomendações foram
implementadas corretamente e que os riscos foram mitigados.
\clearpage
\clearpage\chapter{Da Unidade Administrativa de Tecnologia da Informação e
Comunicação (TIC) do
CRP-06}

\section{Do Grupo Gestor da Tecnologia da Informação e
Comunicação}

Visa coordenar e acompanhar as ações de Tecnologia da Informação e
Comunicação do CRP-06, bem como propor estratégias para o melhor uso dos
sistemas, dispositivos e fluxos da autarquia na gestão da informação e
comunicação.

Deve ser composto no mínimo por:

\begin{enumerate}

\item
    diretoria;
\item
    Comissão de Auditoria e Controle Interno --- CACI;
\item
    gerências;
\item
    coordenação da Unidade Administrativa de Tecnologia da Informação.
\end{enumerate}

\section{Da Unidade Administrativa de Tecnologia da
Informação}

Composta pelas/os trabalhadoras/es efetivas/os da Unidade Administrativa
da Tecnologia da Informação do CRP-06 (Resolução CRP-06 nº~03/2022),
visa otimizar e facilitar o atendimento das/dos usuárias/os, bem como
melhorar a qualidade do serviço prestado por intermédio de formulação e
acompanhamento de indicadores, primando pelo uso de dados abertos e
legíveis por máquina.

\subsection{Contatos para suporte}

\begin{longtable}[]{@{}
>{\raggedright\arraybackslash}p{(\linewidth - 6\tabcolsep) * \real{0.3532}}
>{\raggedright\arraybackslash}p{(\linewidth - 6\tabcolsep) * \real{0.3080}}
>{\centering\arraybackslash}p{(\linewidth - 6\tabcolsep) * \real{0.1848}}
>{\centering\arraybackslash}p{(\linewidth - 6\tabcolsep) * \real{0.1540}}@{}}
\toprule\noalign{}
\begin{minipage}[b]{\linewidth}\centering
\textbf{E-mail}
\end{minipage} & \begin{minipage}[b]{\linewidth}\centering
                 \textbf{Descrição do suporte}
               \end{minipage} & \begin{minipage}[b]{\linewidth}\centering
                                  \textbf{Horário}
                                \end{minipage} & \begin{minipage}[b]{\linewidth}\centering
                                                   \textbf{Responsável}
                                                 \end{minipage}                                               \\
\midrule\noalign{}
\endhead
\bottomrule\noalign{}
\endlastfoot
informatica@crpsp.org.br                  & Rede

Monitoramento

Segurança

Equipamentos

Zimbra

NextCloud                                 & 09:00 às 18:00                            & André

Caio

Carlos                                                                                                                                    \\
sei@crpsp.org.br                          & Sistema SEI!                              & 09:00 às 18:00                            & André

Caio

Carlos                                                                                                                                    \\
sistemas@crpsp.org.br                     & BR Conselhos

BRC Serviços Online                       & 09:00 às 18:00                            & Wellington

Denise

Carlos                                                                                                                                    \\
carlos.vasconcelos@crpsp.org.br           & Novos projetos                            & 10:00 às 19:00                            &
Carlos                                                                                                                                    \\
\end{longtable}

\subsubsection{\texorpdfstring{Suporte à rede (internet e
\emph{wi-fi}), monitoração de segurança, equipamentos e serviços Zimbra
e
Nextcloud}{Suporte à rede (internet e wi-fi), monitoração de segurança, equipamentos e serviços Zimbra e Nextcloud}}
\normalfont

O suporte à rede, monitoração de segurança, equipamentos, serviços
Zimbra, Nextcloud e SEI! será prestado pela equipe de Tecnologia da
Informação, de segunda a sexta-feira, das 9h às 18h.

Chamados deverão ser direcionados para o \emph{e-mail}:
informatica@crpsp.org.br.

\subsubsection{Suporte ao SEI!}

O suporte ao SEI! Será prestado pela equipe de Tecnologia da Informação,
de segunda a sexta-feira, das 9h às 18h.

Chamados deverão ser direcionados para o \emph{e-mail}:
sei@crpsp.org.br.

\subsubsection{Suporte aos sistemas BR Conselhos e BRC Serviços
Online}

O suporte aos serviços ligados ao BRC será prestado pelo trabalhador
Wellington, de segunda a sexta-feira, das 9h às 18h.

Chamados deverão ser direcionados para o \emph{e-mail}:
sistemas@crpsp.org.br.

\subsubsection{Novos projetos}

Todos os novos serviços que forem demandados para a TI que não sejam
chamados, conforme descrito nos itens anteriores, e que demandem
atividades técnicas e de sistema, serão tratados como novos projetos de
TI e deverão obedecer ao seguinte cronograma:

\begin{enumerate}

\item
    a solicitação deverá ser direcionada para o e-mail
    carlos.vasconcelos@crpsp.org.br, com descrição do projeto e o máximo
    possível de informações, como cronograma esperado, informações
    quantitativas, etc.;
\item
    após a abertura do chamado, será organizada uma reunião de kick off
    com todas/os as/os envolvidas/os;
\item
    a reunião deverá apresentar o escopo do projeto, definir o cronograma,
    a avaliação de riscos e a matriz de responsabilidades;
\item
    após a reunião, o projeto será criado no Deck, no sistema de nuvem
    Nextcloud;
\item
    Todas/os as/os envolvidas/os terão acesso ao projeto no Deck, que irá
    gerenciar todas as atividades, responsabilidades, cronograma,
    comunicação e entrega do projeto.
\end{enumerate}

\paragraph{Dos resultados
esperados}

\begin{enumerate}

\item
    Visibilidade das etapas do projeto por todas/os as/os envolvidas/os;
\item
    escopo, documentos e cronograma controlados por uma única ferramenta;
\item
    identificação dos incidentes que poderão impactar no projeto;
\item
    visibilidade dos projetos em andamento, o que permitirá elaborar
    cronogramas mais consistentes;
\item
    acompanhamento dos riscos identificados no planejamento, o que irá
    trazer maior segurança para o andamento do projeto.
\end{enumerate}

\clearpage\chapter{Das normas de controle de
acesso}

\begin{enumerate}

\item
    Os \emph{logins} e senhas da conta institucional são pessoais e
    protegem a identidade da/o usuária/o, e não poderão ser compartilhados
    com outras pessoas em nenhuma hipótese. Tais dispositivos evitam e
    previnem que uma pessoa se faça passar por outra perante o CRP-06;
\item
    não é permitida a criação de contas\emph{/}listas de transmissão ou de
    distribuição em nome de coletivos ou unidades administrativas que não
    possam identificar a pessoa física que receberá ou responderá
    pessoalmente pela conta;
\item
    o uso dos \emph{logins} e/ou senhas de outra pessoa constitui infração
    criminal nos moldes da legislação vigente;
\item
    a/o usuária/o vinculado a tais dispositivos será responsável pelo seu
    uso correto perante a autarquia e a legislação;
\item
    é expressamente proibido o compartilhamento de \emph{login}/senha para
    funções de administração de sistemas.
\item
    a Unidade Administrativa de Tecnologia da informação responde pela
    criação da identidade lógica das/os usuárias/os na instituição.
\item
    ao realizar o primeiro acesso no ambiente do CRP-06 e no \emph{e-mail}
    institucional, a/o usuária/o deverá trocar imediatamente a sua senha,
    conforme as orientações apresentadas;
\item
    todos os acessos devem ser imediatamente bloqueados quando se tornarem
    desnecessários ou em caso de desligamento de qualquer conselheira/o,
    trabalhadora/trabalhador e/ou colaboradora/colaborador;
\item
    eventualmente, os dados da/o usuária/o desligada/o poderão ser
    acessados, somente com a autorização da diretoria, a título de
    \emph{backup} das informações;
\item
    a Unidade de Gestão de Pessoas, a Unidade de Compras e/ou a Unidade de
    Secretaria deverão imediatamente comunicar o desligamento à Unidade de
    TI, a fim de que essa providência seja tomada;
\item
    em caso de esquecimento de sua senha, a/o usuária/o deverá solicitar
    formalmente uma nova junto à Unidade de TIC.
\end{enumerate}

\clearpage\chapter{Dos recursos tecnológicos (computadores, notebooks,
celulares,
tablets)}

\begin{enumerate}

\item
    Os equipamentos do CRP-06 disponibilizados às/aos usuárias/os devem
    ser utilizados exclusivamente para as atividades institucionais;
\end{enumerate}

\begin{enumerate}

\setcounter{enumi}{38}
\item
    é proibido todo procedimento de manutenção física ou lógica,
    instalação, desinstalação, configuração ou modificação de
    \emph{hardware} e/ou \emph{software}, sem o conhecimento prévio e o
    acompanhamento de uma/um técnica/o da TI do CRP-06;
\item
    as áreas que necessitarem de manutenção deverão solicitá-la
    previamente por meio da abertura do chamado técnico à TI por
    \emph{e-mail};
\item
    todas as atualizações e correções de segurança do sistema operacional,
    \emph{softwares} ou aplicativos, somente poderão ser realizadas após a
    devida validação no respectivo ambiente de homologação, cabendo à
    equipe de TI a disposição destas atualizações em ambiente de produção;
\item
    os sistemas e computadores devem ter versões apenas do \emph{software}
    antivírus institucional instaladas, ativadas e atualizadas. Em caso de
    suspeita de vírus ou problemas na funcionalidade, a/o usuária/o deverá
    acionar o departamento técnico responsável mediante registro de
    chamado;
\item
    não será responsabilidade do CRP-06 a violação de acessos pessoais
    e/ou não compatíveis com as atividades desempenhadas pela/o usuária/o,
    como compras pessoais \emph{on-line}, acesso a \emph{sites} de banco
    etc.;
\item
    a transferência e/ou a divulgação de qualquer \emph{software},
    programa ou instruções de computador para terceiros, por qualquer meio
    de transporte (físico ou lógico), somente poderá ser realizada com a
    devida identificação da/o solicitante, via Unidade de TI;
\item
    arquivos pessoais e/ou não correspondentes as atividades
    institucionais --- fotos pessoais, músicas e vídeos, entre outros ---
    não deverão ser copiados/movidos para os \emph{drives} de rede;
\item
    caso identificada a existência desses arquivos, eles poderão ser
    eventualmente excluídos mediante comunicação prévia à/ao usuária/o;
\item
    os documentos institucionais deverão ser salvos em volumes de rede. De
    forma suplementar, poderá ser utilizado o serviço Nextcloud da conta
    institucional da/o usuária/o para realizar cópia das informações
    previamente disponibilizadas em rede;

    \begin{enumerate}
      \def\labelenumii{\roman{enumii}.}
      \item
            tais arquivos, se gravados apenas nos computadores --- por exemplo,
            no \emph{drive} C: ou em outro meio físico ou virtual não validado
            pela TI --- não terão garantia de \emph{backup} e poderão ser
            perdidos caso ocorra uma falha no computador. Neste caso, a
            responsabilidade é da/o própria/o usuária/o;
    \end{enumerate}
\end{enumerate}

\begin{enumerate}

\item
    as/os usuárias/os de contas privilegiadas (conselheiros e gestores)
    não devem, sem a prévia solicitação e a autorização da coordenação de
    TI, executar nenhum tipo de comando ou programa que venha a
    sobrecarregar os serviços existentes na rede institucional. Em caso de
    dúvida, a TI deverá ser acionada;
\item
    no uso dos computadores, equipamentos e recursos de informática,
    algumas regras devem ser necessariamente atendidas:
\end{enumerate}

\begin{enumerate}
\def\labelenumi{\roman{enumi}.}
\item
    A utilização de senha de segurança para acesso os dispositivos;
\item
    todos os computadores de uso individual deverão ter senha na BIOS para
    restringir o acesso de colaboradoras/es não autorizados. Tais senhas
    serão definidas pela Tecnologia da informação, que terá acesso a estas
    para manutenção dos equipamentos;
\item
    as/os usuárias/os devem informar à Unidade de TI qualquer
    identificação de dispositivo estranho conectado ao seu computador;
\item
    é vedada a abertura ou o manuseio de computadores ou outros
    equipamentos de informática para qualquer tipo de reparo que não seja
    realizado por uma/um técnica/o da TI ou por terceiros devidamente
    contratados para o serviço;
\item
    uma vez assumida a responsabilidade como usuária/o de informações, a
    pessoa deverá manter a configuração do equipamento disponibilizado,
    seguindo os devidos controles de segurança exigidos pela Política de
    Segurança da Informação e pelas normas específicas da autarquia;
\item
    quando não estiverem sendo utilizados, todos os terminais de
    computador, \emph{notebook} e impressoras deverão ser protegidos por
    senha;
\item
    todos os recursos tecnológicos adquiridos pelo CRP-06 devem ter,
    imediatamente, suas senhas padrões (\emph{default}) alteradas pela TI;
\end{enumerate}

\begin{enumerate}

\setcounter{enumi}{47}
\item
\item
\item
    fica proibido o uso de computadores e recursos tecnológicos do CRP-06
    para:
\end{enumerate}

\begin{enumerate}
\def\labelenumi{\roman{enumi}.}
\item
    tentar obter acesso não autorizado a outro computador, servidor ou
    rede;
\end{enumerate}

\begin{enumerate}
\def\labelenumi{\roman{enumi}.}
\setcounter{enumi}{7}
\item
    burlar quaisquer sistemas de segurança;
\item
    acessar informações confidenciais sem explícita autorização do
    proprietário;
\item
    vigiar secretamente conteúdo indevido por meio de \emph{softwares}
    analisadores de pacotes;
\item
    interromper um serviço, servidores ou rede de computadores por meio de
    qualquer método ilícito ou não autorizado;
\item
    usar qualquer tipo de recurso tecnológico para cometer ou ser cúmplice
    de atos de violação, propaganda política, assédio sexual ou moral,
    perturbação, caos institucional, difamação, truculência, manipulação,
    ou supressão de direitos autorais ou propriedades intelectuais;
\item
    hospedar material de conteúdo pornográfico, racista, machista,
    capacitista ou que viole a moral, a ordem pública, e ainda, que viole
    a dignidade humana e as leis nacionais ou tratados internacionais de
    que o Brasil seja signatário;
\item
    utilizar \emph{software} pirata, atividade considerada delituosa de
    acordo com a legislação nacional;
\item
    transmitir mensagens em massa sem endereçar os interessados no
    conteúdo.
\end{enumerate}

As infrações terão como encaminhamento as medidas punitivas adotadas
previstas.

\clearpage\chapter{\texorpdfstring{Do correio eletrônico ---
\emph{e-mail}}{Do correio eletrônico --- e-mail}}

\begin{enumerate}

\item
    O uso do correio eletrônico é para fins exclusivamente institucionais
    relativos à tramitação de informações, e usará o domínio
    @crpsp.org.br;
\item
    o \emph{e-mail} deve ser um meio rápido de resposta, e não um sistema
    de armazenamento de dados, que devem ser armazenados na rede;
\item
    é proibida a utilização do \emph{e-mail} institucional para
    deliberações de qualquer natureza, tramitação de documentos
    confidenciais ou sigilosos e mensagens com fins expositivos,
    truculentos, caluniosos ou difamatórios contra qualquer pessoa;
\end{enumerate}

\begin{enumerate}
\def\labelenumi{\roman{enumi}.}
\item
    listas de distribuição\emph{/}transmissão somente serão usadas para
    tramitação de comunicação institucional de memorandos e demais
    documentos oficiais, exclusivamente de Unidades Administrativas e
    entre trabalhadoras/es, sendo vedado seu uso para solicitação de
    informações, para deliberações de qualquer natureza, tramitação de
    documentos confidenciais ou sigilosos e mensagens com fins
    expositivos, truculentos, caluniosos ou difamatórios;
\end{enumerate}

\begin{enumerate}

\item
\item
\item
\item
    em caso de desligamento, o serviço será ser bloqueado, cabendo a
    notificação prévia desta ação à Secretaria\emph{/}Gestão de Pessoas,
    que comunicará à TI;
\item
    cabe à TI realizar a gestão das contas sob domínio institucional,
    sendo efetuado bloqueio por três meses no caso de não utilização do
    \emph{e-mail}, licença e/ou desligamento;
\end{enumerate}

\begin{enumerate}
\def\labelenumi{\roman{enumi}.}
\item
    os dados deverão ser solicitados pela/o gestora/gestor imediata/o
    antes do término deste prazo. A não manifestação dentro do prazo pode
    tornar a informação irrecuperável;
\end{enumerate}

\begin{enumerate}

\setcounter{enumi}{5}
\item
    é VEDADO o uso do correio eletrônico para:
\end{enumerate}

\begin{enumerate}
\def\labelenumi{\roman{enumi}.}
\item
    enviar mensagem em nome e/ou conta de outra/o usuária/o sem a devida
    autorização;
\item
    armazenar documentos institucionais que deveriam constar da rede de
    computadores;
\item
    deliberações de qualquer natureza, tramitação de documentos
    confidenciais ou sigilosos e mensagens com fins expositivos,
    truculentos, caluniosos ou difamatórios;
\item
    enviar qualquer mensagem por meios eletrônicos que torne seu remetente
    e/ou o CRP-06 e suas unidades vulneráveis às ações civis ou criminais;
\item
    divulgar informações não autorizadas ou imagens de tela, sistemas,
    documentos e afins, sem autorização expressa e formal concedida pela/o
    proprietária/o desse ativo de informação;
\item
    falsificar informações de endereçamento, adulterar cabeçalhos para
    esconder a identidade de remetentes e/ou destinatários, com o objetivo
    de evitar as punições previstas;
\item
    produzir, transmitir ou divulgar mensagem que:
\end{enumerate}

\begin{enumerate}
\def\labelenumi{\alph{enumi}.}
\item
    contenha qualquer ato ou forneça orientação que conflite ou contrarie
    os interesses públicos do CRP-06.
\item
    contenha ameaças eletrônicas, como: \emph{spam}, \emph{mail bombing},
    vírus de computador, \emph{phishing};
\item
    contenha arquivos com código executável (.exe, .com, .bat, .pif, .js,
    .vbs, .hta, .src, .cpl, .reg, .dll, .inf) ou qualquer outra extensão
    que represente um risco à segurança;
\item
    vise obter acesso não autorizado a outro computador, servidor ou rede;
\item
    vise interromper um serviço, servidores ou rede de computadores por
    meio de qualquer método ilícito ou não autorizado;
\item
    vise burlar qualquer sistema de segurança;
\item
    vise vigiar secretamente ou assediar outra/o usuária/o;
\item
    vise acessar informações confidenciais sem explícita autorização da/o
    proprietária/o;
\item
    vise acessar indevidamente informações que possam causar prejuízos a
    qualquer pessoa;
\item
    inclua imagens criptografadas ou de qualquer forma mascaradas;
\item
    tenha conteúdo considerado impróprio, obsceno ou ilegal, seja de
    caráter calunioso, difamatório, truculento, ofensivo, violento,
    pornográfico, que gere caos institucional, entre outros;
\item
    contenha perseguição preconceituosa baseada em gênero, raça,
    incapacidade física ou mental ou outras situações protegidas;
\item
    tenha fins políticos locais ou nacionais;
\item
    inclua material protegido por direitos autorais sem a permissão da/o
    detentora/detentor dos direitos;

    \begin{enumerate}
      \def\labelenumii{\roman{enumii}.}
      \item
      \item
      \item
      \item
      \item
      \item
      \item
            a notificação de infrações e a aplicação de eventuais medidas
            punitivas caberão às gerências e à diretoria, em conjunto com a
            equipe de Gestão de Pessoas;
      \item
            nos casos mais graves, será aberta sindicância para apurar os fatos
            e tomar as medidas legais cabíveis, sem prejuízo de instauração de
            processo administrativo disciplinar, caso os fatos sejam evidentes e
            não necessitem de apuração prévia;
      \item
            as mensagens de correio eletrônico deverão obrigatoriamente incluir
            assinatura, com cargo e/ou função.
    \end{enumerate}
\end{enumerate}

\clearpage\chapter{Do uso da internet}

Qualquer informação que é acessada, transmitida, recebida ou produzida
na internet está sujeita a divulgação e auditoria de órgão competente.
Portanto, o CRP-06, em total conformidade com a legislação vigente,
reserva-se o direito de monitorar, quando necessário, e, sempre que
possível, registrar todos os \emph{links} de acessos dentro de seu
domínio institucional.

Os equipamentos, tecnologia e serviços fornecidos para o acesso à
internet são de propriedade da instituição, que pode analisar e, se
necessário, bloquear qualquer arquivo, \emph{site}, correio eletrônico,
domínio ou aplicação armazenados na rede/internet, estejam eles em disco
local, na estação ou em áreas privadas da rede.

Quanto ao uso propriamente dito da rede, fica definido que:

\begin{enumerate}

\item
    toda tentativa de alteração dos parâmetros de segurança, por qualquer
    usuária/o, sem o devido credenciamento e a autorização para tal, será
    julgada inadequada e os riscos relacionados serão informados à/ao
    colaboradora/colaborador e à/ao respectiva/o gestora/gestor;
\end{enumerate}

\begin{enumerate}

\setcounter{enumi}{50}
\item
    o uso de qualquer recurso para atividades ilícitas poderá acarretar
    ações administrativas e as penalidades decorrentes de processos civil
    e criminal. Nesses casos, a instituição cooperará ativamente com as
    autoridades competentes;
\item
    a internet disponibilizada pela instituição às/aos suas/seus
    colaboradoras/es, independentemente de sua relação contratual, não
    pode ser utilizada para fins pessoais;
\item
    é proibida a divulgação e/ou o compartilhamento indevido de
    informações da área administrativa em listas de discussão, redes
    sociais, \emph{sites}, salas de bate-papo, comunicadores instantâneos
    ou qualquer outra tecnologia correlata que venha a surgir na internet;
\item
    as/os usuárias/os com acesso à internet somente poderão fazer o
    \emph{download} de conteúdos ligados diretamente às suas atividades no
    CRP-06. Para isso, deverão providenciar o que for necessário para
    regularizar a licença e o registro de eventuais programas com a
    autorização a supervisão local;
\item
    esses \emph{softwares} deverão passar previamente pelo crivo de
    segurança da equipe de TI. O uso, a instalação, a cópia ou a
    distribuição não autorizada de \emph{softwares} que tenham direitos
    autorais, marca registrada ou patente na internet são expressamente
    proibidos;
\item
    qualquer \emph{software} não autorizado baixado poderá ser excluído
    pela TI;
\item
    as/os usuárias/os não poderão, em hipótese alguma, utilizar os
    recursos do CRP-06 para fazer o \emph{download} ou distribuição
    de \emph{software} ou dados pirateados;
\item
    o \emph{download} e a utilização de programas de entretenimento, jogos
    ou músicas (em qualquer formato) poderão ser realizados por
    usuárias/os que tenham atividades profissionais relacionadas a essas
    categorias. Para tal, serão criados grupos de segurança, cujas/os
    integrantes deverão ser definidos pelas/os respectivas/os gestoras/es,
    a fim de viabilizar esse acesso especial;
\item
    como regra geral, material de cunho sexual para pesquisa acadêmica não
    poderá ser exposto, armazenado, distribuído, editado, impresso ou
    gravado por meio de qualquer recurso;
\end{enumerate}

\begin{enumerate}
\def\labelenumi{\roman{enumi}.}
\item
    caso seja necessário, grupos de segurança deverão ser criados para
    viabilizar esse perfil de usuária/o especial e seus integrantes
    definidos pelas/os respectivas/os gestoras/es e informados ao setor de
    Tecnologia da Informação;
\end{enumerate}

\begin{enumerate}

\setcounter{enumi}{59}
\item
    usuárias/os com acesso à internet não poderão efetuar \emph{upload}
    e/ou \emph{download} de qualquer \emph{software} licenciado ao CRP-06
    ou de dados de sua propriedade a seus/suas parceiras/os e clientes;
\item
    os serviços de comunicação instantânea (Skype, Meeting, Hangout
    pessoal e afins) serão inicialmente disponibilizados às/os usuárias/os
    e poderão ser bloqueados caso a/o gestora/gestor requisite formalmente
    à coordenação de TI, uma vez que há ferramentas corporativas para
    comunicação disponíveis para todas as contas de \emph{e-mail} (Zoom e
    InMail);
\item
    não é permitido acesso a \emph{sites} de \emph{proxy}. Quando
    necessário, solicitar à Unidade de Tecnologia da Informação;
\item
    toda e qualquer pesquisa institucional promovida eletronicamente, por
    meio de formulários, deverá ser previamente notificada à Unidade
    Jurídica, para que seja submetida a análise de conformidade com a LGPD
    (Lei Geral de Proteção de Dados) e, desta forma, ser verificado o
    atendimento às questões legais.
\end{enumerate}

\clearpage
\chapter{Dos dispositivos móveis
institucionais}

O CRP~SP deseja facilitar a mobilidade e o fluxo de informação entre
suas/seus usuárias/os. Por isso, permite a utilização de equipamentos
portáteis, exclusivamente para a finalidade profissional.

Por ``dispositivo móvel'', entende-se qualquer equipamento eletrônico
com atribuições de mobilidade de propriedade da instituição, como:
\emph{notebooks}, \emph{smartphones}, \emph{tablets}, cartões de
memória, HDs externos e \emph{pen drives}.

O CRP~SP, na qualidade de proprietário dos equipamentos e/ou serviços
fornecidos, reserva-se o direito, caso seja necessário, de
inspecioná-los a qualquer momento e realizar manutenção de segurança.

Portanto, a/o usuária/o assume o compromisso de não utilizar, revelar ou
divulgar a terceiros, de modo algum, direta ou indiretamente, em
proveito próprio ou de terceiros, qualquer informação, confidencial ou
não, que tenha ou venha a ter conhecimento em razão de suas funções na
instituição, mesmo depois de terminado o vínculo institucional.

Ficam definidas, ainda, as seguintes diretrizes:

\begin{enumerate}

\item
    toda/o usuária/o deverá realizar periodicamente cópia de segurança
    (\emph{backup}) dos dados de seu dispositivo móvel institucional,
    armazenando os dados em rede;
\item
    não é permitida a guarda de informações institucionais em dispositivos
    pessoais. Em caso de dúvidas quanto a este processo, entre em contato
    com a equipe de TI;
\item
    a TI só assegurará \emph{backup} dos dados dispostos em rede
    institucional. Toda/o usuária/o deverá utilizar senhas de bloqueio
    automático para seu dispositivo móvel;
\item
    não será permitida, em nenhuma hipótese, a alteração da configuração
    dos sistemas operacionais dos equipamentos;
\item
    a/o usuária/o será responsabilizada/o por manter ou utilizar quaisquer
    programas e/ou aplicativos que não tenham sido instalados ou
    autorizados por um técnico da TI ou que estejam em desacordo com a
    política institucional do CRP~SP;
\item
    a reprodução não autorizada dos \emph{softwares} instalados nos
    dispositivos móveis fornecidos pela instituição constituirá uso
    indevido do equipamento e infração legal aos direitos autorais do
    fabricante;
\item
    é permitido o uso de rede \emph{wi-fi} de locais conhecidos pela/o
    usuária/o, como sua casa, hotéis, fornecedores e clientes;
\item
    não é recomendado o compartilhamento de dados em redes \emph{wi-fi}
    desconhecidas;
\item
    é responsabilidade da/o usuária/o notificar imediatamente sua/seu
    gestora/gestor direta/o e o departamento de Tecnologia da Informação
    do furto ou roubo de um dispositivo móvel fornecido pelo CRP~SP. Além
    disso, deverá procurar a ajuda das autoridades policiais, registrando
    um boletim de ocorrência;
\item
    a/o usuária/o deverá estar ciente de que o uso indevido do dispositivo
    móvel caracterizará a assunção de todos os riscos pela má utilização e
    sua responsabilização por danos dolosos, diretos ou indiretos,
    presentes ou futuros, que venha a causar ao CRP~SP e/ou a terceiros;
\item
    a/o usuária/o que desejar utilizar equipamentos portáteis particulares
    ou adquirir acessórios e posteriormente conectá-los à rede do CRP~SP
    deverá submeter previamente tais equipamentos ao processo de
    autorização de sua/seu superiora/superior imediata/o para liberação
    junto ao setor de Tecnologia da Informação.
\end{enumerate}

\section{Da telefonia móvel}

Temos à disposição o serviço de telefonia móvel institucional para uso
de usuárias/os previamente autorizadas/os pela diretoria. O serviço é
restrito às atribuições funcionais.

O \emph{e-mail} institucional da/o usuária/o será o identificador
necessário para validação do aparelho.

A Unidade de TIC não se responsabiliza pela utilização de aplicativos
e/ou conteúdo não condizente com os valores institucionais ou
atribuições profissionais.

Em caso de avarias, furto, roubo ou perda do dispositivo, bem como
\emph{backup} e/ou garantia de dados, a/o colaboradora/colaborador é a/o
responsável pela abertura do boletim de ocorrência e aviso ao
departamento de TI, conforme termo de adesão, para que as medidas de
segurança possam ser adotadas. Por isso, é recomendado que a/o
colaboradora/colaborador peça periodicamente a realização de
\emph{backup} completo dos dados do aparelho em nuvem ou em
\emph{desktop}/\emph{notebook}.

\clearpage\chapter{Do backup}

\begin{enumerate}

\item
    O \emph{backup} de dados dos sistemas institucionais, rede, servidores
    e \emph{e-mail} é de responsabilidade do setor de Tecnologia da
    Informação;
\item
    todos os \emph{backup}s devem ser automatizados por sistemas de
    agendamento que sejam preferencialmente executados fora do horário
    comercial, nas chamadas ``janelas de \emph{backup}'' (diariamente
    entre as 19h e as 5h). Esses são períodos de nenhum ou pouco acesso de
    usuárias/os e de processos automatizados aos sistemas de informática;
\item
    as antigas mídias de \emph{backup} (como DAT, LTO, DVD, CD e outros)
    devem ser acondicionadas em local seco, climatizado, seguro e dentro
    do \emph{datacenter/}Cedoc;
\item
    com o objetivo de excluir mídias que possam apresentar riscos de
    gravação ou de restauração decorrentes do uso prolongado, além do
    prazo recomendado pelo fabricante, o tempo de vida e uso das mídias de
    \emph{backup} deve ser monitorado e controlado pelas/os responsáveis e
    de acordo com as orientações do fabricante;
\item
    mídias que apresentam erros devem primeiramente ser formatadas e
    testadas. Caso o erro persista, deverão ser inutilizadas;
\item
    os \emph{backup}s imprescindíveis, críticos, para o bom funcionamento
    do CRP~SP exigem uma regra de retenção especial, conforme recomendação
    da ISO 27000, seguindo, assim, as determinações fiscais e legais
    existentes no país;
\item
    na situação de erro de \emph{backup} ou restauração, é necessário que
    ele seja feito logo no primeiro horário disponível, assim que a/o
    responsável tenha identificado e solucionado o problema e respeitar a
    publicação de acordo de nível de serviço para a entrega da informação;
\item
    a restauração de dados, quando solicitada, poderá se dar em ambiente
    de teste e/ou mídia externa para não comprometer demais informações
    atualizadas;
\item
    demandas fora deste padrão deverão ocorrem em mídias físicas
    alternativas e de forma suplementar mediante a solicitação prévia de
    gestora/gestor.
\item
    não cabe à TI a responsabilidade por \emph{backup}, formatação ou
    edição de dados computacionais disponíveis em outras formas de
    armazenamento que não os volumes da rede institucional.
\end{enumerate}

\clearpage\chapter{Das disposições finais}

Questões não previstas nesta Política deverão ser tratadas
individualmente e pela Unidade de TIC do CRP-06 e encaminhados para
deliberação da diretoria.

Há necessidade de revisões semestrais deste documento, a fim de mantê-lo
atualizado às necessidades institucionais.

Assim como a ética, a segurança deve ser entendida como parte
fundamental da cultura interna do CRP-06, ou seja, qualquer incidente de
segurança será subentendido como ação contra o bem público e a
responsabilidade regida pela autarquia.

\backmatter
\chapter{Do armazenamento em nuvem --- Nextcloud}

O Nextcloud é a ferramenta institucional de código aberto para
sincronização de arquivos, por meio do qual são compartilhados
\emph{softwares} para todas/os as/os usuárias/os, e fornece uma solução
de compartilhamento e sincronização de arquivos segura e em conformidade
com os servidores do CRP-06.

Serve para compartilhar um ou mais arquivos e pastas em seu computador e
sincronizá-los com seu servidor Nextcloud, e será acessado pelo domínio
nuvem.crpsp.org.br.

Os dados de \emph{login} e senha serão ofertados pela Unidade de TI para
as/os usuárias/os autorizados.

É proibido o uso de outras formas de armazenamento em nuvem que não
sejam por contas contratadas pelo CRP-06.

A interface da/o usuária/o Nextcloud contém os seguintes campos e
funções:

\begin{enumerate}

\item
    \textbf{menu de seleção de aplicativos}: localizado no canto superior
    esquerdo, nele se encontram todos os aplicativos disponíveis em sua
    instância do Nextcloud. Clicar em um ícone de aplicativo irá
    redirecioná-lo para o aplicativo;
\item
    \textbf{campo ``informações do aplicativo''}: localizado na barra
    lateral esquerda, fornece filtros e tarefas associadas ao aplicativo
    selecionado. Por exemplo, ao usar o aplicativo Arquivos, pode-se
    acessar um conjunto especial de filtros para localizar rapidamente
    seus arquivos, como aqueles que foram compartilhados com você e
    aqueles que você compartilhou com outras pessoas. Outros aplicativos
    terão itens diferentes;
\item
    \textbf{visualização do aplicativo}: campo central principal na
    interface da/o usuária/o Nextcloud, que exibe o conteúdo ou os
    recursos da/o usuária/o para o aplicativo selecionado;
\item
    \textbf{barra de navegação}: localizada sobre a janela de visualização
    principal (a visualização do aplicativo), esta barra fornece um tipo
    de navegação estrutural que permite que você migre para níveis mais
    altos da hierarquia de pastas até o nível raiz (home);
\item
    \textbf{botão ``novo''}: localizado na barra de navegação, o botão
    ``novo'' permite criar novos arquivos, novas pastas ou fazer upload de
    arquivos;
\item
    \textbf{campo ``pesquisar''}: clique na lupa no canto superior direito
    para pesquisar arquivos e entradas do aplicativo atual;
\item
    \textbf{menu ``contatos''}: fornece uma visão geral de seus contatos e
    das/os usuárias/os em seu servidor. Dependendo dos detalhes fornecidos
    e dos aplicativos disponíveis, pode ser usado para iniciar uma chamada
    de vídeo diretamente com elas/eles ou enviar e-mails;
\item
    \textbf{botão visualização de grade}: botão formado por quatro
    pequenos quadrados que alterna a visualização da grade para pastas e
    arquivos;
\item
    \textbf{menu ``configurações''}: clique na imagem do seu perfil,
    localizada à direita do campo ``pesquisar'', para abrir o menu
    suspenso ``configurações''. A página de configurações fornece os
    seguintes recursos:
\end{enumerate}

\begin{enumerate}
\def\labelenumi{\roman{enumi}.}
\item
    \emph{links} para baixar aplicativos de \emph{desktop} e móveis;
\end{enumerate}

\begin{enumerate}
\def\labelenumi{\roman{enumi}.}
\setcounter{enumi}{15}
\item
    uso do servidor e espaço disponível;
\item
    gerenciamento de senha;
\item
    configurações de nome, \emph{e-mail} e foto do perfil;
\item
    gerenciar navegadores e dispositivos conectados;
\item
    membras/os do grupo;
\item
    configurações de linguagem da interface;
\item
    gerenciar notificações;
\item
    ID da nuvem federada e botões de compartilhamento de mídia social;
\item
    gerenciador de certificados SSL/TLS para armazenamento externo;
\item
    suas configurações de autenticação com dois fatores;
\item
    informação da versão do Nextcloud.
\end{enumerate}

Os arquivos Nextcloud, com a interface da \emph{web} Nextcloud, servem
para criar, visualizar, editar, excluir, compartilhar e recompartilhar
arquivos.

A/o administradora/administrador do Nextcloud tem a opção de desativar
esses recursos. Portanto, se algum deles estiver faltando no sistema,
pergunte à/ao administradora/administrador do servidor.

Podem-se atribuir \emph{tags} aos arquivos. Para criar \emph{tags}, abra
um arquivo na visualização ``detalhes''. Em seguida, digite suas
\emph{tags}. Para inserir mais de uma etiqueta, pressione a tecla
``enter'' depois de criar cada etiqueta. Todas as \emph{tags} são
\emph{tags} do sistema e são compartilhadas por todas/os as/os
usuárias/os no servidor Nextcloud.

\section{Visualizar arquivos}

É permitido exibir arquivos de texto não compactados, arquivos
OpenDocument, vídeos e arquivos de imagem nos visualizadores
incorporados do Nextcloud, clicando-se no nome do arquivo. Pode haver
outros tipos de arquivo que você pode visualizar se o administrador do
Nextcloud os tiver ativado. Se o Nextcloud não puder exibir um arquivo,
ele inicia um processo de \emph{download} e baixa o arquivo no seu
computador.

\section{Navegando dentro do
Nextcloud}

Navegar pelas pastas no Nextcloud é tão simples quanto clicar em uma
pasta para abri-la e usar o botão voltar no seu navegador para passar
para o nível anterior. O Nextcloud também fornece uma barra de navegação
na parte superior do campo ``arquivos'' para uma navegação rápida.

\section{Ícones de status do
compartilhamento}

Qualquer pasta compartilhada é marcada com o ícone de sobreposição
``compartilhado''. Os compartilhamentos de \emph{links} públicos são
marcados com um elo de corrente. As pastas não compartilhadas não são
marcadas.

\section{Selecionando arquivos ou
pastas}

Pode ser selecionado um ou mais arquivos ou pastas clicando-se nas
caixas de seleção. Para selecionar todos os arquivos no diretório atual,
clique na caixa de seleção localizada na parte superior da lista de
arquivos.

Ao selecionar vários arquivos, você pode excluir todos eles ou baixá-los
como um arquivo ZIP, usando os botões ``excluir'' ou ``baixar'' que
aparecem na parte superior.

\section{Filtrando a visualização de
arquivos}

A barra lateral esquerda na página ``arquivos'' contém vários filtros
para classificar e gerenciar rapidamente seus arquivos.

\section{Todos os arquivos}

A visualização padrão; exibe todos os arquivos aos quais você tem
acesso.

\section{Favoritos}

Arquivos ou pastas marcados com a estrela amarela.

\section{Compartilhado com você}

Exibe todos os arquivos compartilhados com você por outra/o usuária/o ou
grupo.

\section{Compartilhado com outras
pessoas}

Exibe todos os arquivos que você compartilhou com outras/os usuárias/os
ou grupos.

\section{Compartilhado por link}

Exibe todos os arquivos compartilhados por você por meio de \emph{link}
público.

\section{Movendo arquivos}

Mover arquivos e pastas arrastando e soltando-os em qualquer diretório.

\clearpage\chapter{Dos produtos Google}

O Google Workspace é um conjunto de soluções de nuvem que contribui para
a melhor operação dos processos de trabalho da autarquia.

Promove a melhoria na produtividade e na colaboração, integrando
inúmeras soluções disponíveis em qualquer lugar e em qualquer
dispositivo.

Oferece proteção das informações da forma mais segura possível, passando
por auditorias e cumprindo com os elevados padrões de segurança do
setor.

A plataforma Google só será utilizada para gerar \emph{links} de
reuniões, em função da acessibilidade, e para fins de procedimentos da
Comissão de Ética, Comissão de Orientação e Fiscalização, diretoria e
secretaria, com a função de armazenar as gravações previstas nas
normativas do CFP.

As contas contratadas ficarão sobre a responsabilidade das unidades
acima citadas, não sendo permitido o compartilhamento de senhas e
arquivos considerados confidenciais, como as gravações.

Oferece as seguintes ferramentas:

\section{Docs}

Para criar e editar documentos de texto, planilhas, tabelas e outros de
forma \emph{on-line} e \emph{off-line}. Armazenamento seguro
sincronizado com o Google Drive.

\section{Drive}

Para guardar documentos, fotos e planilhas de forma segura com o serviço
de armazenamento e sincronização de arquivos.

\section{Meet}

Para videoconferências com segurança e acessibilidade.

\section{Forms}

Para coleta de pesquisas e informações por meio de formulários criados
para pessoas ou grupos.

\section{Sheets}

Ferramenta para criação e edição de planilhas, gráficos e tabelas.
Serviço sincronizado com Google Docs.

\section{Agenda}

Para criação de agendas compartilhadas que permitem ver quando outras
pessoas estão disponíveis\,e programar reuniões com convites automáticos
por \emph{e-mail}.

\clearpage\chapter{Dos produtos Zimbra/InMail}

A plataforma Zimbra Mail é o domínio oficial e preferencial do CRP-06 e
garante um \emph{e-mail} seguro, aplicação de mensagens instantâneas
individuais ou para grupos, realização de videoconferência com
possibilidade de gravação e guarda das informações do \emph{chat},
criação e compartilhamento de calendário e agenda, controle de tarefas,
compartilhamento de documentos e vinculação com o armazenamento em
nuvem.

Estes recursos estão disponíveis para plataformas de \emph{webmail} e
\emph{desktop} e como aplicativos para dispositivos móveis (celulares e
\emph{tablets}).

Todos os recursos são de uso exclusivo para atividades institucionais,
não sendo permitida a utilização e o compartilhamento para fins
pessoais.

\section{\emph{E-mail}}

O e\emph{-mail} oferece a possibilidade de selecionar uma data e horário
específicos para envio, por meio do recurso ``enviar mais tarde'' que
fica localizado no \emph{menu} suspenso.

Além da comunicação por \emph{e-mail}, sua caixa de entrada do Zimbra é
um ponto de partida para compartilhamento e colaboração de conteúdo,
bate-papos e videoconferências. O cliente \emph{web} responsivo do
Zimbra pode ser acessado de qualquer lugar e em qualquer dispositivo, e
permite uma transição perfeita entre os aplicativos do Zimbra. As/os
usuárias/os podem redigir vários \emph{e-mail}s e alternar facilmente
entre conversas com colegas de trabalho, edição de documentos ou revisão
do calendário.

Se você se esqueceu de adicionar algo ao seu \emph{e-mail} logo após
clicar no botão ``enviar'' ou simplesmente se esqueceu de incluir um
destinatário importante, o recurso de ``desfazer envio'' do Zimbra
oferece alguns segundos para que você possa interromper o envio da
mensagem para fazer quaisquer alterações. Os \emph{e-mail}s Zimbra
também podem ser programados para envio em data e hora futuras.

Também é possível recuperar itens excluídos que foram para a lixeira,
por até 30 dias. Porém, esse recurso tem que ser configurado pela/o
usuária/o.

\section{Chat integrado com dispositivo de mensagem instantânea
---
ImMail}

O ImMail é o dispositivo oficial e exclusivo do CRP-06 para troca de
mensagens instantâneas para fins de avisos, recados e compartilhamento
de \emph{links} para reuniões.

É vedado o uso do ImMail para fins não institucionais ou passíveis de
causar ações acusatórias, truculentas, difamatórios ou que produzam caos
institucional.

\section{Lista de tarefas}

Para criar listas de tarefas já com detalhes de informações, por
exemplo: datas de vencimento, prioridades, entre outros pontos.

\section{Calendário}

Para marcar uma reunião \emph{on-line} abrindo o calendário e, por meio
do ``assistente de agendamento'', verificar os dias e horários
disponíveis de outras/os colaboradoras/es da sua empresa.

É possível utilizar atalhos de cópia, por meio do botão direito. Essa
função permite um ``copia e cola'' dos compromissos (com as/os
respectivas/os convidadas/os, anexos etc.).

\section{Pasta para armazenamento em
nuvem}

Anexar arquivos direto do ``Porta de Arquivos''.

Salvar anexos de \emph{e-mail}, também, no ``Porta de Arquivos''.

Carregar arquivos para acessá-los de forma instantânea no Zimbra Mail.

\section{Pesquisa}

Possibilita pesquisar o que se quiser na conta Zimbra (inclusive
anexos), aproveitar os \emph{menu}s robustos para filtrar resultados de
pesquisas, salvar cada pesquisa feita, se necessário, e ver os
resultados (que não desaparecem quando você clica fora da janela).

\section{Outros recursos (em quesitos mais técnicos e
administrativos)}

Implantar o \emph{e-mail}, na nuvem, no local ou, ainda, como sistema
híbrido.

Fazer o gerenciamento de recursos da/o usuária/o final, além de cotas,
políticas de armazenamento, por meio de classe de serviço (CoS).

Ter proteção anti-\emph{spam} e antivírus integrados.

Além disso, o e\emph{-mail} Zimbra suporta IMAP/POP, CaIDAV e CardDAV. A
plataforma Zimbra também possui o Admin Console, que permite a
administração da plataforma de qualquer lugar.

\section{Controle e customização
total}

O e\emph{-mail} Zimbra permite o gerenciamento técnico em nível de
infraestrutura. O console de administração permite fazer o gerenciamento
dos seus \emph{e-mail}s. Também oferece a interface de linha de comando,
a conhecida CLI.

O Zimbra possui anti-\emph{spam} e antivírus integrados na ferramenta.
Defina suas próprias políticas de acesso, anti-\emph{spam} e demais
configurações.

Com o gerenciamento é possível criar restrição de acesso aos
\emph{e-mail}s por horário, conforme você preferir.

O Zimbra permite controle total do sistema, podendo ter sua instalação
customizada de acordo com a necessidade do seu negócio. Com o servidor
dedicado, é possível mudar qualquer regra --- como acesso, política de
senha, controle anti-\emph{spam} e muitas outras.

\section{Acesso mobile e desktop}

Para uma melhor experiência da/o usuária/o, é possível acessar seu
ambiente Zimbra nas plataformas \emph{desktop} e \emph{mobile}.

\section{Segurança}

Segurança de dados e informações é outro ponto crucial para uma
plataforma de \emph{e-mail}. Nesse caso, você também pode ficar
tranquilo, porque o Zimbra Mail possui recursos de alta segurança.

\section{Estabilidade}

A estabilidade pode ser comprovada em qualquer dispositivo móvel.

\section{Compartilhamento de
calendário}

No ambiente colaborativo do Zimbra Mail, é possível compartilhar
calendários. Dessa forma, todas/os as/os profissionais podem ver os
horários disponíveis dos outros.

É possível sugerir outros horários clicando com o botão direito, no
evento, para copiar e colar todas as informações (inclusive anexos)
referentes a ele.

\section{Interface intuitiva}

Permite que as/os usuárias/os possam mudar facilmente de navegadores ---
além de \emph{desktop}, celular e \emph{tablet} --- enquanto aproveitam
toda a experiência de colaboração e \emph{e-mail} estável.

Personaliza a interface da/o usuária/o com cores, fontes, logotipos e
outros detalhes de sua preferência.

\clearpage\chapter{Dos sistemas}

\section{BRConselhos --- BRC}

Trata-se do sistema oficial do Sistema Conselhos de Psicologia. É uma
solução integrada para gestão dos conselhos, que permite o controle de
serviços gerenciais, financeiros e de fiscalização de forma remota, com
garantia de segurança e integridade de acessos.

Apresenta as seguintes funcionalidades:

\begin{enumerate}

\item
    \textbf{requerimentos on-line}: as/os psicólogas/os inscritas/os no
    CRP~SP têm à disposição um sistema de autoatendimento on-line que
    permite solicitar primeira inscrição ou inscrições secundárias, fazer
    atualização cadastral, solicitar declarações, cancelar inscrições,
    reativar inscrições suspensas ou canceladas, emitir certidão
    profissional, entre outras funcionalidades;
\item
    \textbf{atendimento}: o atendimento é a porta de entrada da/do
    psicóloga/o ou da pessoa jurídica para o CRP~SP, oferecendo acesso aos
    serviços disponíveis, como CRP Acolhe, inscrição de pessoa física ou
    jurídica, renegociação de débitos, cancelamento, reativação, entre
    outros. Permite geração de relatórios;
\item
    \textbf{COE}: possibilita iniciar e acompanhar todas as fases dos
    processos éticos, mantendo os registros, fases processuais, documentos
    e resultados. Geração de relatórios, por processo, fase processual,
    registro de audiências, andamentos, mediação e pessoas mediadoras,
    entre outros;
\item
    \textbf{COF}: viabiliza o registro de controle das demandas de
    fiscalização, manutenção dos cadastros de temas, andamentos, fases
    processuais, questionários, documentos e relatórios;
\item
    \textbf{financeiro}: registro das anuidades, controle dos pagamentos
    via integração bancária, registro e disponibilização da posição
    financeira de cada uma das inscrições, processos de cobrança,
    manutenção da estrutura financeira, relatórios de cobrança, anuidades
    geradas, controle dos boletos, repasses ao CFP, Inscrição em Dívida
    Ativa (termo de inscrição e emissão de CDA --- Certidão de Dívida
    Ativa);
\item
    \textbf{jurídico}: manutenção da dívida ativa, gestão e controle de
    processos, inclusão e exclusão de inadimplentes, andamento dos
    processos, relatórios dos processos, geração dos documentos e
    processos de cobrança;
\item
    \textbf{contabilidade}: manutenção da integração contábil, gerando os
    relatórios necessários para o processo de registro contábil dos
    resultados, como entradas efetivas, receitas por conta (mapa
    contábil), fundo de seção, fechamento contábil/financeiro.
\end{enumerate}

\section{Sistema Eletrônico de Informação
(SEI!)}

O sistema de gerenciamento de processos SEI! pode ser acessado por meio
dos principais navegadores e permite o acesso remoto.

O sistema pretende garantir a eficiência dos processos sistematizados
entre órgãos municipais, estaduais e federais; a variedade de formatos e
tamanhos de documentos compatíveis; o controle de nível de acesso; e a
tramitação simultânea de processos em múltiplas unidades. Por meio de
uma interface intuitiva, a plataforma oferece ainda funcionalidades
específicas, como controle de prazos, estatísticas da unidade, tempo do
processo, pesquisa de inteiro teor, acompanhamento especial,
textos-padrão, assinatura em bloco, organização de processos em bloco, e
outros. Foi instituído pelo Decreto nº~8.539/2015.

DECRETO Nº~8.539, DE 8 DE OUTUBRO DE 2015 --- Dispõe sobre o uso do meio
eletrônico para a realização do processo administrativo no âmbito dos
órgãos e das entidades da administração pública federal direta,
autárquica e fundacional.

No âmbito do Executivo Federal, o extinto MPOG coordenou a implementação
do Processo Eletrônico Nacional (PEN), com a adoção do Sistema
Eletrônico de Informação (SEI!) pelos ministérios e outros órgãos
vinculados ao Executivo federal até 2020.

O Processo Eletrônico Nacional (PEN) foi uma iniciativa conjunta de
órgãos e entidades de diversas esferas da administração pública, com o
intuito de construir uma infraestrutura pública de processos e
documentos administrativos eletrônicos, objetivando a melhoria no
desempenho dos processos do setor público, com ganhos em agilidade,
produtividade, transparência, satisfação da/o usuária/o e redução de
custos. Existiu até 2020. Introduziu práticas inovadoras no setor
público: eliminação do uso de papel como suporte físico para documentos
institucionais e disponibilização de informações em tempo real. Foi
composto por algumas grandes ações, sendo o Sistema Eletrônico de
Informações (SEI), desenvolvido pelo Tribunal Regional Federal da 4ª
Região, a principal entrega.

O SEI! é uma plataforma que engloba um conjunto de módulos e
funcionalidades que promovem a eficiência administrativa. A solução é
cedida gratuitamente para instituições públicas e permite transferir a
gestão de documentos e de processos eletrônicos administrativos para um
mesmo ambiente virtual.

\subsection{Vantagens}

\begin{enumerate}

\item
    \textbf{portabilidade}: 100\% \emph{web} e pode ser acessado por meio
    dos principais navegadores do mercado;
\item
    \textbf{acesso remoto}: pode ser acessado remotamente por diversos
    tipos de equipamentos, como microcomputadores, \emph{notebooks},
    \emph{tablets} e \emph{smartphones} de vários sistemas operacionais.
    Isso possibilita que as/os usuárias/os trabalhem à distância;
\item
    \textbf{acesso de usuárias/os externos}: gerencia o acesso de
    usuárias/os externos, permitindo que tomem conhecimento dos documentos
    e, por exemplo, assinem remotamente contratos e outros tipos de
    processos;
\item
    \textbf{controle de nível de acesso}: gerencia a criação e o trâmite
    de processos e documentos restritos e sigilosos, conferindo o acesso
    somente às unidades envolvidas ou a usuárias/os específicos;
\item
    \textbf{tramitação em múltiplas unidades}: incorpora novo conceito de
    processo eletrônico, que rompe com a tradicional tramitação linear,
    inerente à limitação física do papel. Várias unidades podem ser
    demandadas, tomar providências e manifestar-se simultaneamente;
\item
    \textbf{funcionalidades específicas}: controle de prazos, ouvidoria,
    estatísticas da unidade, tempo do processo, base de conhecimento,
    pesquisa em todo teor, acompanhamento especial, inspeção
    administrativa, modelos de documentos, textos-padrão, sobrestamento de
    processos, assinatura em bloco, organização de processos em bloco,
    acesso externo etc.;
\item
    \textbf{sistema intuitivo}: estruturado com boa navegabilidade e
    usabilidade. \emph{Software} público gratuito. Além disso, o acordo de
    cooperação técnica (ACT) assinado entre o CFP e o MPOG e a estrutura
    de TI já instalada ou em aquisição pelo CFP, conforme deliberação em
    Assembleia das Políticas, da Administração e das Finanças (Apaf),
    permite a adesão gratuita dos conselhos regionais ao PEN e ao SEI!
    para uso do sistema, mediante manifestação de interesse. Acordo já
    formalizado e utilizando uma estrutura única compartilhada;
\item
    \textbf{integração ao Sistema Conselhos de Psicologia}:
\end{enumerate}

\begin{enumerate}
\def\labelenumi{\roman{enumi}.}
\item
    Elaboração de minutas de documentos e tramitações poderá ser executada
    por funcionárias/os do CRP a pedido de conselheiras/os
    (preferencialmente por \emph{e-mail}, para registro);
\item
    \emph{e-mail}s com respostas de conselheiras/os poderão ser anexados
    aos processos por funcionárias/os;
\item
    conselheiras/os poderão assinar documentos elaborados por
    funcionárias/os e colocados em blocos internos, mesmo remotamente;
\item
    reuniões e plenárias poderão ter suas pautas, com subsídios,
    organizadas e apresentadas integralmente via SEI!;
\item
    processos de diversas áreas/comissões poderão ser elaborados e
    tramitar entre os setores, com registro das deliberações e dos
    encaminhamentos executados por cada ator envolvido.
\end{enumerate}

\section{Sistema de passagens e diárias
(Sispad)}

O Sispad permite efetuar o controle das solicitações de viagens,
autorizações das/os gestoras/es, emissão das passagens aéreas, pagamento
das diárias e demais despesas com viagens custeadas pelo Conselho para
conselheiras/os, diretoras/es, colaboradoras/es e convidadas/os.

O controle do sistema abrange todas as fases do processo e está
integrado aos demais módulos desenvolvidos para as áreas financeira,
administrativa e contábil.

A/O conselheira/conselheiro, colaboradora/colaborador ou
trabalhadora/‌trabalhador deverá ingressar no Sispad com seu \emph{login}
e senha e seguir o Fluxo de Solicitação de Ressarcimento e Pagamento e
da Prestação de Contas ao CRP-06.

\end{document}
